\chapter{The porting method}\label{ch:port}

The port moves about two hundred loop nests of WRF onto the GPU, spread over some fifty source files, and every one
of them must give the same bits as before. No single test can find a mistake in a change of that size. The method
therefore breaks the work into steps so small that each can be verified on its own, and it surrounds every step with
checks that run without human judgement. This chapter describes the method; \cref{ch:subsys} applies it to WRF's
subsystems.

\section{Two views of one source}

Every ported file compiles in two ways:
\begin{itemize}
  \item the \emph{CPU view}, with the macro \code{WRF_GPU} undefined, is what CPU-REF compiles;
  \item the \emph{GPU view}, with \code{WRF_GPU} defined, is what GPU-REPRO compiles.
\end{itemize}
The rule is that the CPU view must compile to exactly what the unported source compiled to, statement for statement.
Everything the GPU needs that is not a directive goes under \texttt{\#ifdef WRF\_GPU}. Directive lines are comments for a
compiler without \code{-acc}, so a kernel directive in front of an unchanged loop needs no \texttt{\#ifdef}.

The reference text is fixed: the \emph{CPU-view base commit} (\code{port/agent/cpu_view_base}). The static checker
\code{arith_guard} preprocesses each file without \code{WRF_GPU}, compares every executable statement with the same
file at the base commit, and fails on any difference. A small list of additions is allowed: \code{USE} of the port's
modules, \code{CALL gpu_*} and \code{CALL bt_*}, island includes, and route tests that do no arithmetic.

Some changes must reach both views: replacing the intrinsic \code{EXP} by \code{rp_exp}, or fixing an undefined
result. These \emph{shared refactors} are allowed only with evidence: the refactored CPU build must equal the
unrefactored one bit for bit on the test windows (tests T-SHARED-1 to T-SHARED-9). Only then does the base commit move,
by a recorded procedure.

\section{Residency: the state lives on the device}

The model state of the case takes about 31~GB (d01 6.8~GB, d02 24.7~GB including the fire grid; the estimator's
numbers in \code{port/RESULTS.md}). Copying it over PCIe would take about a second, while a d02 step must take
milliseconds. The port therefore keeps the whole state on the device for the whole run:
\begin{itemize}
  \item \textbf{Allocation.} WRF's Registry generates the code that allocates every field of a domain
    (\code{tools/gen_allocs.c}, \cref{ch:arch}). The port extends the generator to add, after each allocation, an
    \code{enter data} that creates the device copy and fills it. The allocation of a domain thus also allocates its
    device copy, for every field of every package in use.
  \item \textbf{Generated update lists.} A second generator writes, for each purpose, the list of \code{update}
    statements over the fields concerned: all fields, the lateral boundary arrays, the fields the nest forcing reads
    from the parent.
  \item \textbf{Tables.} Module data that kernels read (the Noah parameter tables, RRTMG's absorption coefficients,
    WSM6's constants, the fire flags) is declared on the device and uploaded once after initialization.
\end{itemize}

\paragraph{Sync points.} The host still runs initialization, I/O and part of the nest forcing. Data moves only at
named points (\code{plan.md}, P1.5):
\begin{center}
\small
\begin{tabular}{lll}
\toprule
Point & Where & What moves \\
\midrule
S1 & after the initial state is read & all fields and tables to the device \\
S2 & when a nest starts & parent to the host; both domains to the device \\
S3 & first statement of \code{med_hist_out} & the fields of that history stream to the host \\
S4 & first statement of \code{med_restart_out} & the fields of the restart stream to the host \\
S5 & after a lateral boundary read & the boundary arrays to the device \\
S6 & around the nest forcing & what the host forcing reads and writes \\
\bottomrule
\end{tabular}
\end{center}
The history and restart transfers walk the domain's list of state variables at run time, because the content of a
stream can change through \code{iofields_filename} (\cref{ch:iobuild}).

\paragraph{The world flag.} While the port is incomplete, most routines still run on the host. One flag,
\code{gpu_world_host}, records where the current values of the model arrays are. During Phases 1 to 4 the model
lives on the host: \code{solve_em} downloads the state at its start and uploads it at its end (the \emph{bracket}),
and each ported routine moves its own arguments to the device and back. In Phase 5, when every routine of the step is
ported, the bracket is removed and the flag flips: the model lives on the device, and only a routine switched off
moves data.

\section{Routes and islands}

Every ported routine has a \emph{route}, an integer \code{R_<NAME>} in \code{frame/module_gpu_route.F} (101 routes
for this case). Every kernel of the routine carries \code{if(gpu_on(R_<NAME>))}, and the environment variables
\code{WRF_GPU_OFF} and \code{WRF_GPU_ONLY} switch routes at run time. A route that is off runs the same GPU-view code
on the host, through OpenACC's host fallback.

A routine that runs where its data are not needs an \emph{island}: copies of its arguments to where it runs, and
back. The islands are generated from each routine's argument list (\code{port/tools/gen_island.py}), never written by
hand. \Cref{alg:island} shows what they do.

\begin{algorithm}[ht]
\caption{The island of a routine with route $R$ (generated; GPU view only).}
\label{alg:island}
\begin{algorithmic}[1]
\State $\mathit{isl} \gets$ (\code{gpu_world_host} and \code{gpu_on}($R$)) or (not \code{gpu_world_host} and not \code{gpu_on}($R$))
\If{$\mathit{isl}$} \Comment{the routine runs where its data are not}
  \State copy every array argument to where the routine runs
  \State \code{gpu_world_host} $\gets$ not \code{gpu_world_host} \Comment{nested routines see the right world}
\EndIf
\State the routine's statements and kernels
\If{$\mathit{isl}$} \Comment{before every \code{RETURN} and at the end}
  \State \code{gpu_world_host} $\gets$ not \code{gpu_world_host}
  \State copy every argument that is not \code{INTENT(IN)} back
\EndIf
\end{algorithmic}
\end{algorithm}

Islands make every routine portable alone. A routine can be ported, tested and committed while its neighbours still
run on the host, and two routines can be ported in parallel by different people. The price is speed: each island
moves its arguments twice per call. Islands are scaffolding, and they move nothing once the whole step runs on the
device.

\begin{keypoint}
Between the entry and the exit of an island, no host code may read or write a moved array. Its host copy is stale
while the routine runs on the device, and a host write would be overwritten by the exit copy. This mistake is
invisible to every CPU check, because host and ``device'' memory are the same there; only a GPU run reveals it.
\end{keypoint}

\section{Checking one routine}

Three tests check a ported routine, from the narrowest to the widest.

\paragraph{The harness.} \code{port/h100/harness.sh} calls one routine alone with random inputs, three times: the
CPU view, the GPU view on the host, the GPU view on the device. The outputs must agree bit for bit. It takes seconds
and reaches code paths that a short model run may not.

\paragraph{The call check.} Inside a real model run, the island of a routine can check its own calls
(\code{frame/module_gpu_callcheck.F}). \Cref{alg:callcheck} shows how. Because both executions start from the same
inputs, any difference is the kernel's own: a race, a missing \code{private}, a wrong index range, a reordered
statement, or data the island does not move.

\begin{algorithm}[ht]
\caption{The call check of route $R$ (\code{WRF_GPU_CALLCHECK}=$R$:$n$).}
\label{alg:callcheck}
\begin{algorithmic}[1]
\State save every argument the routine may change
\State run the routine with $R$ on the device; keep its results
\State restore the saved arguments
\State run the routine again with every route on the host, as CPU-REF would
\State compare every argument bit for bit; report the first differing element with both values in hexadecimal
\end{algorithmic}
\end{algorithm}

\paragraph{T-AB and T-TRACE.} Over a window of the dev case, T-AB compares the run with the route on against the run
with the route off: the routine's device execution against its host execution, on real data, over many calls and
both domains. T-TRACE then compares GPU-REPRO with CPU-REF over the same window. The comparison uses the bit-hash
traces of \cref{ch:repro}; the first differing line names the step, the RK stage, the checkpoint and the field.

T-AB and T-TRACE test different things. T-AB compares the GPU view with itself and finds mistakes in kernels and
islands. T-TRACE compares the GPU view with the original source and finds mistakes in the restructuring itself, such
as a column kernel whose statements are in a different order than the slab loop's.

\section{Templates}

A loop nest is ported by recognizing its pattern and applying the matching template
(\code{port/agent/CODING_STANDARD.md}, section~5). Each template has a tested example in \code{port/tests}, and most
have \emph{mutants}: deliberately wrong versions that the test must reject.

\subsection{A: pointwise}

Every iteration writes only its own point and reads nothing that another iteration writes. The directive goes in
front of the loop and nothing else changes (\cref{lst:calc-alt-acc}). When the source has several nests in a row, each
becomes its own kernel, in source order; the end of a kernel is the synchronization that the order needs. 92 of the
174 rows of the port's kernel table are of this kind.

\subsection{B: rolling buffer}

Advection computes the flux through each cell face and takes its difference. The CPU code does this row by row, and
keeps the fluxes of the previous row in a two-row buffer, swapping two indices (\cref{alg:advect-y}). Iteration $j$
uses what iteration $j-1$ computed, so the rows cannot run in parallel. Template B removes the dependence by storing
all faces:
\begin{itemize}
  \item kernel Y1 computes the flux of every face $j$ into a 3D work array \code{fqy3(i,k,j)}, with the same branch
    for the same $j$ (5th order inside, 3rd or 2nd order next to the boundaries) and the same expressions;
  \item kernel Y2 computes, for every row, the tendency from \code{fqy3(i,k,j) - fqy3(i,k,j-1)}.
\end{itemize}
Every face flux is still computed once, from the same operands, and every tendency is the same difference. Only the
storage changed. The price is one 3D work array, written once and read twice.

\begin{figure}[ht]
\begin{lstlisting}[basicstyle=\ttfamily\footnotesize]
   ! K-ADVU-Y2: flux divergence of rows j_start .. j_end (the source's j-1)
!$acc parallel loop gang vector collapse(3) if(dev) default(none) present(tendency, fqy3, msfux) &
!$acc& firstprivate(i_start, i_end, j_start, j_end, kts, ktf, rdy) private(mrdy)
   DO j = j_start+1, j_end+1
   DO k=kts,ktf
   DO i = i_start, i_end
      mrdy=msfux(i,j-1)*rdy   ! ADT eqn 44, 2nd term on RHS
      tendency(i,k,j-1) = tendency(i,k,j-1) - mrdy*(fqy3(i,k,j)-fqy3(i,k,j-1))
   ENDDO
   ENDDO
   ENDDO
\end{lstlisting}
\caption{The second kernel of Template B for the $y$-flux of \code{advect_u} (from \code{port/tests/templates/t_tmpl_b.F90}).
The statements are the source's; only the buffer index became the row index.}
\label{lst:tmpl-b}
\end{figure}

The $x$-fluxes have no such buffer: the two faces of a cell are computed in the same iteration. A thread can compute
both faces of its own cell; each face is then computed twice, by two neighbouring threads, with the same operations
on the same operands, which gives the same bits.

\subsection{C: column}

A vertical recurrence, such as the tridiagonal sweep of the acoustic solver (\cref{alg:tridiag}) or the
integration of $\omega$ in \code{calc_ww_cp}, cannot be parallel in $k$. The source usually runs it over a slab:
for each $j$, an array \code{s(its:ite,kts:kte)} is filled level by level for all $i$ at once. Template C gives each
$(i,j)$ column one thread, and turns the slab into a private column:
\begin{itemize}
  \item a slab \code{s(its:ite,kts:kte)} becomes a private fixed-size array \code{s_col(WRF_KMAX+1)};
  \item a per-row vector \code{v(its:ite)} becomes a private scalar;
  \item initializations done once before the $j$ loop move into the column;
  \item statement groups that the source runs over different $i$ ranges keep their own range, as an \code{IF} on $i$;
  \item the $k$ loops run sequentially inside the thread, in the source's order and direction (\texttt{!\$acc loop seq}).
\end{itemize}
The private arrays have a fixed size because a private array of run-time size would live in the slow device heap.
\code{WRF_KMAX} is 64, and the startup check stops a run with more levels.

\Cref{lst:tmpl-c} shows the start of the column kernel of \code{calc_coef_w}, which sets up the coefficients of the
implicit $w$ equation. In the source, \code{cof(i)} is a vector over a row; here it is the private scalar \code{cofs}.

\begin{figure}[ht]
\begin{lstlisting}[basicstyle=\ttfamily\scriptsize]
!$acc parallel loop gang vector collapse(2) if(gpu_on(R_CALC_COEF_W)) default(none) &
!$acc& present(a,alpha,gamma,mut,c1h,c2h,c1f,c2f,cqw,rdn,rdnw,c2a) &
!$acc& firstprivate(i_start,i_end,j_start,j_end,kde,dts,g,epssm,lid_flag) private(k,kk,b,c,cofs)
      DO j = j_start, j_end
      DO i = i_start, i_end
        k = kde-1
        cofs  = (.5*dts*g*(1.+epssm))**2
        a(i, 2 ,j) = 0.
        ...
        !$acc loop seq
        DO k=2,kde-1
          b = 1.+cqw(i,k,j)*cofs*rdn(k)*( ... )
          c =   -cqw(i,k,j)*cofs*rdn(k)*rdnw(k  )*c2a(i,k,j  )/( ... )
          alpha(i,k,j) = 1./(b-a(i,k,j)*gamma(i,k-1,j))
          gamma(i,k,j) = c*alpha(i,k,j)
        ENDDO
\end{lstlisting}
\caption{Template C: the column kernel of \code{calc_coef_w} (excerpt of \code{port/tests/tools/calc_coef_w_gpu.inc};
the elided expressions are the source's, verbatim).}
\label{lst:tmpl-c}
\end{figure}

\subsection{D: calls inside kernels}

When a loop body calls a procedure, such as \code{fire_ros} for every fire node or a column scheme for every column,
the callee and everything it calls get \texttt{!\$acc routine seq}. The compiler builds a device version that runs in
the calling thread, beside the host version. A device routine may not do I/O, allocate memory, stop the model, or
touch \code{grid} or \code{config_flags}; error branches set an integer code instead, and the host acts on it after
the kernel with the source's own message.

\subsection{G: boundary strips}

Boundary conditions copy values into rows and columns at the edges of the domain (\cref{ch:bdy}). Each strip of the
source becomes its own kernel, in the source order. The order matters: the $y$ strips read the corners that the $x$
strips wrote, so strips of different phases are never merged into one kernel.

\subsection{L: split update}

The positive-definite limiter of \code{advect_scalar_pd} (\cref{alg:pd}) loops over cells, and a cell that would be
emptied scales the fluxes through its outflow faces. A face is shared by two cells. In a parallel loop, one thread
would write the face while its neighbour reads it to test the flux's sign: a data race. It happens to be harmless
here, but only through an argument about signs and zeros that no compiler checks. Template L splits the update in
two:
\begin{itemize}
  \item a cell kernel computes each cell's scale factor, or $-1$ when the cell is not limited;
  \item a face kernel per direction looks at the sign of each face flux, finds the donor cell, and multiplies by the
    donor's factor if the donor lies in the source's cell range and is limited.
\end{itemize}
Each face is then written by exactly one thread, with one multiplication by the same factor as in the source, and only
the faces the source scales are touched. The unit test T-PDLIM compares the original loop and the split version on
random and real states, including zero fluxes and the edges of the range.

\subsection{R: row-parallel}

The ozone interpolation \code{ozn_p_int} searches, for each level, the two pressure levels of the climatology that
bracket it. The source starts the search of column $i$ where the search of column $i-1$ ended. The iterations over
$i$ therefore share state. Template R keeps the source's $i$ loop sequential inside one thread per row $j$, so that
each search starts where the source's started and finds the same bracket.

\subsection{CP: column physics}

A physics scheme is called by a wrapper that loops over rows, copies a slab of the model fields into the scheme's
arrays, calls the scheme for the row, and copies the results back. Template CP (\code{plan.md} section~8.0) replaces
the wrapper's row loop by one kernel over $(j,i)$:
\begin{enumerate}
  \item each thread copies its column into private fixed-size arrays, with the wrapper's own expressions (for WSM6,
    \code{t = th*pii});
  \item it calls the scheme with a single column, \code{its = ite = 1}, \code{kts = 1}, \code{kte = nz};
  \item it copies the results back with the wrapper's expressions (\code{th = t/pii}).
\end{enumerate}
The scheme and its callees become device routines; their automatic arrays sized by the tile become fixed-size, and
every whole-array statement on them (\code{a(:) = 0.}) is rewritten to the active extent \code{a(1:nz)}. The original
wrapper stays in the CPU view, unchanged. \Cref{ch:colphys} discusses the performance of this pattern.

\section{Memory for temporaries}

WRF's routines declare large temporary arrays as \emph{automatic arrays}: arrays sized by the arguments, allocated on
the stack at each call. A kernel cannot use them: they are host memory, created afresh at every call. The port
replaces them in two ways.
\begin{itemize}
  \item \textbf{The scratch pool.} The \code{i1} arrays of \code{solve_em}, which the Registry declares as automatic
    tendency arrays, become pointers into one large array allocated once (\code{frame/module_gpu_scratch.F}). The
    generator that declares them emits the pointer form.
  \item \textbf{Work arrays.} Other large automatic arrays (the flux arrays of advection, the temporaries of the fire
    and radiation drivers) become named arrays allocated once at initialization for the largest domain
    (\code{frame/module_gpu_work.F}). The routine remaps a pointer onto the work array with the bounds the automatic
    array had.
\end{itemize}
Both are compiled into both builds (\code{-DWRF_POOL}) and zero-filled once. So CPU-REF and GPU-REPRO see the same
memory in every temporary, even if some code read a temporary before writing it. The test T-UNINIT shows that no code
of the case does: a build that fills every local with signaling NaNs and a build that fills them with zeros give the
same bits.

\section{Static guards}

Some mistakes can be found by reading the code mechanically. \code{port/gates/static.sh} runs these checks before
every commit:
\begin{description}
  \item[\code{arith_guard}] The CPU view of every file equals the base commit, apart from the allowed additions.
  \item[\code{kernel_lint}] Every kernel has a route, \code{default(none)}, no data clause, a perfect nest under
    \code{collapse}, \code{loop seq} on every inner loop that is not an early-exit search, no I/O, and no intrinsic
    transcendental function (it must be \code{rp_*}); every file with kernels has the island of each route.
  \item[\code{check_verbatim}] The statements of a restructured kernel are the source's statements, character for
    character, apart from index changes the template allows.
  \item[Protection] The tests, gates and comparison tools are checksummed, so that no one, human or agent, can make a
    failing test pass by editing it.
\end{description}

\section{The routine ladder}

Each routine climbs eight rungs (\code{roadmap/plan-4.6/README.md}):
\begin{center}
\small
\begin{tabular}{llll}
\toprule
Rung & What & Machine & Passes when \\
\midrule
L1 & write the GPU view & CLOUD & \code{static.sh} passes \\
L2 & compile both gfortran views & CLOUD & no errors \\
L3 & GPU view equals CPU view on the host & CLOUD & identical bits \\
L4 & compile with \code{nvfortran} & WS-A100 & kernels offloaded, no implicit copies \\
L5 & harness: host, device, CPU view & WS-A100 & three-way identical \\
L6 & T-AB on a dev-case window & WS-A100 & identical traces \\
L7 & T-TRACE against CPU-REF & WS-A100 & bitwise \\
L8 & record in the kernel database & WS-A100 & rows present \\
\bottomrule
\end{tabular}
\end{center}
The first three rungs need no GPU and run in a cloud environment; the others need a GPU and the case data. A routine
is done at rung L7. The machines are named by what they can do:
\begin{itemize}
  \item \textbf{CLOUD}, a cloud coding environment with gfortran and no GPU;
  \item \textbf{WS-A100}, a workstation with two A100 40~GB and 56 host cores: every device test, and every
    CPU-REF reference run on its host cores, so that CPU and GPU results are compared on one machine with one
    compiler;
  \item \textbf{CI}, GitHub Actions, which builds this book.
\end{itemize}
Nothing is developed or tested on the cluster where the case was first run; it supplies the input files, and its
original run is compared with CPU-REF statistically, never bit for bit. The full case does not fit one 40~GB GPU, so
the acceptance run uses the case \code{eaton_mid} (\cref{ch:case}) and the full case follows on both GPUs
(\cref{ch:mgpu}).

\section{Organization}

The plan (\code{roadmap/plan-4.6/}) arranges the work in nine stages and 140 phases:
\begin{itemize}
  \item 0, ground truth: the reproducible CPU reference, the probes, the dev and acceptance cases, and tiers of
    small generated gating cases that both builds run and must match bit for bit, from ideal cases of seconds to
    small real-data cases of minutes (\cref{app:tests});
  \item 1, infrastructure: state, pool, work arrays, tables and sync points, with all compute still on the host;
  \item 2, dynamics; 3, physics; 4, fire;
  \item 5, nest forcing on the device, the end of the bracket, and the full 17-hour acceptance run;
  \item 6, performance without changing a bit (Part~VII of this book);
  \item 7, other fires and a release; 8, the rest of WRF~4.6.0, including several GPUs.
\end{itemize}
Each stage ends with a gate that someone else can repeat. Inside a stage the phases are independent, thanks to the
islands, and can run in parallel.

\paragraph{Code-only runs.} Rungs L1 to L3 need no GPU, so the writing of kernels can be spread over many workers in
parallel. The routines of Phases 1 to 3 were divided into 44 work packages, each owning a fixed set of routines, so
that two workers never edit the same routine. A scope checker enforces the ownership. The first such run, with AI
coding agents as workers, wrote code for 19 of the 44 packages; 234 kernels passed \code{kernel_lint}, and
\code{arith_guard} found no new arithmetic in GPU code (\code{roadmap/plan-4.6/README.md}).

\paragraph{What the first run taught.} The review of that run found problems that no static check could see:
\begin{itemize}
  \item an infrastructure module that the run's own CPU-view changes depended on had not been written, so the CPU view
    compiled but crashed;
  \item kernels written against code paths the smoke case does not take, which stopped on its open boundaries;
  \item incremental builds that kept stale objects after an include file changed, because \code{make} does not track
    includes (\cref{ch:iobuild}).
\end{itemize}
Each became a phase of Stage~0 of the plan. The lesson is that the static checks bound what can go wrong in the
arithmetic, but only builds and runs show whether the pieces fit.
